\chapter*{Preface}
\addcontentsline{toc}{chapter}{Preface}

Machine learning addresses the issue of analyzing, reproducing and predicting various mechanisms and processes observable through experiments and data acquisition. With the impetus of large technological companies in need of leveraging information included in the gigantic datasets that they produced or obtained through user data, with the development of new data acquisition techniques in biology, physics or astronomy, with the improvement of storage capacity and high-performance computing, this field has experienced an explosive growth over the past decades, in terms of scientific production and technological impact.



While it is being recognized in some places as a scientific discipline in itself, machine learning (which has received a few almost synonymic denominations across time, including artificial intelligence, machine intelligence or statistical learning), can also be seen as an interdisciplinary field interfacing techniques from traditional domains such as computer science, applied mathematics, and statistics. From statistics, and more specially nonparametric statistics, it borrows its main formalism, asymptotic results and generalization bounds. It also builds on many classical methods that have been developed for estimation and prediction. From computer science, it involves the construction and implementation of efficient algorithms, programming design and architecture. Finally, machine learning leverages classical methods from linear algebra and functional analysis, as well as from convex and nonlinear optimization, fields within which it had also provided new problems and discoveries. It forms a significant part of the larger field commonly called ``data science,'' which includes methods for storing, sharing and managing data, the development powerful computer architectures for increasingly demanding algorithms, and, importantly, the definition of ethical limits and processes through which data should be used in the modern world.
 
 
This book, which originates from lecture notes of a series of graduate course taught in the Department of Applied Mathematics and Statistics at Johns Hopkins University, adopts a  viewpoint (or bias) mainly focused on the mathematical and statistical aspects of the subject. Its goal is to introduce the mathematical foundations and techniques that lead to the development and analysis of many of the algorithms that are used today. It is written with the hope to provide the reader with a deeper understanding of the algorithms made available to her in multiple machine learning packages and software, and that she will be able to assess their prerequisites and limitations, and to extend them and develop new algorithms. Note that, while adopting a presentation with a strong mathematical flavor, we will still make explicit the details of many important machine learning algorithms.

Unsurprisingly, the book will be  more accessible to a reader with some background in mathematics and statistics. It assumes familiarity with basic concepts in linear algebra and matrix analysis, in multivariate calculus and in probability and statistics. We tried to place a limit at the use of measure theoretic tools, that are avoided up to a few exceptions, which are be localized and be accompanied with alternative interpretations allowing for a reading at a more elementary level. 

The book starts with an introductory chapter that describes notation used throughout the book and serve at a reminder of basic concepts in calculus, linear algebra and probability. It also introduces some measure theoretic terminology, and can be used as a reading guide for the sections that use these tools. This chapter is followed by two chapters offering background material on matrix analysis and optimization. The latter chapter, which is relatively long, provides necessary references to many algorithms that are used in the book, including stochastic gradient descent, proximal methods, etc. 

\Cref{chap:intro}, which is also introductory, illustrates the bias-variance dilemma in machine learning through the angle of density estimation and motivates \cref{chap:general} in which basic concepts for statistical prediction are provided. \Cref{chap:higher} provides an introduction to reproducing kernel theory and Hilbert space techniques that are used in many places, before tackling, with \cref{chap:lin.reg,chap:lin.class,chap:nearest.neighbors,chap:trees,chap:neural.nets}, the description of various algorithms for supervised statistical learning, including linear methods, support vector machines, decision trees, boosting, or neural networks.

\Cref{chap:mcmc}, which presents sampling methods and an introduction to the theory of Markov chains, starts a series of chapters on generative models, and associated learning algorithms. Graphical models and described in \cref{chap:mrf,chap:inf.mrf,chap:bayes.net}. \Cref{chap:var.bayes} introduces variational methods for models with latent variables, with applications to graphical models in \cref{chap:learning.graphical}. Generative techniques using deep learning are presented in \cref{chap:generative}.

\Cref{chap:clustering,chap:dim.red,chap:manifold.learning} focus on unsupervised learning methods, for clustering, factor analysis and manifold learning. The final chapter of the book is theory-oriented and discusses concentration inequalities and generalization bounds.

%This book, which also provides a list of problems with solutions at the end of each chapter, includes material that can be used to develop a two-semester course on the subject. Each chapter is fairly independent of the others given the basic reference chapters 2--5 and 11. 



%
%
% 
%This will not prevent us from detailing some of the important algorithms, and therefore introducing some notions from optimization theory. We will not assume any prerequisite in advanced linear algebra or optimization  and will try to make the presentation self-contained in these domains. In contrast, we will assume sufficient knowledge in probability theory and multivariate statistics, as provided in textbooks such as \cite{rice2007mathematical,lehmann2006theory,van2000asymptotic}. We will, as most as possible, avoid any reliance on measure-theoretic probability. There will, 
%
%The goal of statistical learning (a term that we will use synonymously to machine learning)  is to infer hidden functional relationships or structure from random data. Probably the most important class of problems (both in terms of publication volume and of impact for applications) is that of prediction, whose goal is to approximate the value of a dependent variable by a function of independent variables. 
% Other important problems aim at finding groups (or clusters) within datasets, reducing their dimensions by including the data in some linear or nonlinear ``manifold'' or providing good visualization tools. Many of this problems are limited by the ``curse of dimensionality'', which expresses the fact that no algorithm can be expected to work uniformly well in large dimensions.